\uuid{57qN}
\titre{Variables discrètes}
\auteur{worms}
\organisation{exo7}
\licence{CC BY-NC-SA 4.0}
\datecreate{2017-08-28}
\niveau{L2}
\module{Probabilité et statistique}
\chapitre{Probabilité discrète}
\sousChapitre{Variable aléatoire discrète}
\difficulte{}
\theme{Variables discrètes}
\exo7qcmid{worms.worms.032}

\contenu{
\begin{question}
\qclassification{260.03}
\qkeeporder
Soit $X$ une variable aléatoire à valeurs dans $\Nn$ vérifiant $\mathbb{P}(X\geq 3)=64\%$. Que peut-on dire de la médiane de $X$ ?
\begin{answers}
\bad{Elle vaut 3.}
\bad{Elle est inférieure ou égale à 3.}
\good{Elle est supérieure ou égale à 3.}
\bad{On manque d'informations pour affirmer l'une des propositions ci-dessus.}
\end{answers}
\begin{explanations}
Elle ne vaut pas forcément $3$ (par exemple, on peut très bien avoir $\mathbb{P}(X=3)=11\%$ donc $\mathbb{P}(X\leq 3)=\mathbb{P}(X\leq 2)+\mathbb{P}(X= 3)=(1-0.64)+0.11=0.47=47\%$ donc $\mathbb{P}(X\leq 3)<1/2$), et elle est $\geq 3$ car si $k\leq 2$, par croissance de la fonction de répartition on a  $\mathbb{P}(X\leq k) \leq \mathbb{P}(X\leq 2)= 1-\mathbb{P}(X>2) =1-\mathbb{P}(X\geq 3) = 1-0.64 < 1/2$.
\end{explanations}
\end{question}
}
